% GENERATED FILE. Edit canonical lessons, then run npm run generate:books.
% canonical-source-hash: fnv1a64:ea7abb87442697f5
% canonical-lessons: GE-C430-wanze, GE-C430-kakerlake, GE-C430-muecke, GE-C430-stechen, GE-C430-krabbeln

\chapter{Ungebetene Gäste}
\label{ch:ge-ungebetene-gaste}
\hlchaptermodality{430}

\begin{chapteropening}
\textbf{By the end of this chapter:} \emph{I can complain about insects in a room: the bedbug, the cockroach, the mosquito, being bitten and crawling.}

Complete the last lesson of chapter 430: I can complain about insects in a room: the bedbug, the cockroach, the mosquito, being bitten and crawling.
\end{chapteropening}

\section[die Wanze]{die Wanze — the bedbug; a hidden microphone}

\label{lesson:GE-C430-wanze}



\begin{quote}
Before the new one: say the German for to squeak, then the German for faulty.
\end{quote}

\subsection*{You'll want to know: die Wanze}
\textbf{die Wanze} — \textquotedblleft{}the bedbug; a hidden microphone\textquotedblright{}.

\begin{quote}
Im Bett waren Wanzen, und er hat schlecht geschlafen.
\end{quote}

\begin{quote}
\emph{There were bedbugs in the bed, and he slept badly.}
\end{quote}

Plural: \emph{die Wanzen}. A hidden microphone is also \emph{eine Wanze}.

\begin{cousinweb}
\emph{Wanze} goes back to Old High German \emph{wantlūs} 'wall louse' (EtymDB), from \emph{Wand} 'wall' and \emph{Laus} 'louse'.
\end{cousinweb}

\begin{grammarlens}[title={What you've built}]
One more word for telling, explaining and describing at B1.
\end{grammarlens}

\subsection*{Your turn}
\begin{itemize}
\raggedright
  \item \emph{Say it:} \emph{die Wanze}
  \item \emph{Say it:} \emph{die Wanze}, once more
  \item \emph{Say it:} a sentence of your own with \emph{die Wanze}
  \item \emph{Recall:} say the German for to squeak, then the German for faulty, then say \emph{die Wanze} again
\end{itemize}

\begin{tcolorbox}[breakable,colback=teal!4,colframe=teal!35!black,title={Before you move on}]
What does die Wanze mean? (\textquotedblleft{}The bedbug; a hidden microphone\textquotedblright{}.) Say it once more.
\end{tcolorbox}
\section[die Kakerlake]{die Kakerlake — the cockroach}

\label{lesson:GE-C430-kakerlake}



\begin{quote}
Before the new one: say the German for faulty, then the German for the bedbug.
\end{quote}

\subsection*{You'll want to know: die Kakerlake}
\textbf{die Kakerlake} — \textquotedblleft{}the cockroach\textquotedblright{}.

\begin{quote}
In der Küche läuft eine Kakerlake über den Boden.
\end{quote}

\begin{quote}
\emph{A cockroach runs across the kitchen floor.}
\end{quote}

Plural: \emph{die Kakerlaken}. Also \emph{die Schabe}.

\begin{cousinweb}
\emph{Kakerlake} probably came into German through Dutch \emph{kakkerlak}, a word sailors brought back from overseas. EtymDB lists it without a source, and the deeper origin is not settled.
\end{cousinweb}

\begin{grammarlens}[title={What you've built}]
One more word for telling, explaining and describing at B1.
\end{grammarlens}

\subsection*{Your turn}
\begin{itemize}
\raggedright
  \item \emph{Say it:} \emph{die Kakerlake}
  \item \emph{Say it:} \emph{die Kakerlake}, once more
  \item \emph{Say it:} a sentence of your own with \emph{die Kakerlake}
  \item \emph{Recall:} say the German for faulty, then the German for the bedbug, then say \emph{die Kakerlake} again
\end{itemize}

\begin{tcolorbox}[breakable,colback=teal!4,colframe=teal!35!black,title={Before you move on}]
What does die Kakerlake mean? (\textquotedblleft{}The cockroach\textquotedblright{}.) Say it once more.
\end{tcolorbox}
\section[die Mücke]{die Mücke — the mosquito, the gnat}

\label{lesson:GE-C430-muecke}



\begin{quote}
Before the new one: say the German for the bedbug, then the German for the cockroach.
\end{quote}

\subsection*{You'll want to know: die Mücke}
\textbf{die Mücke} — \textquotedblleft{}the mosquito, the gnat\textquotedblright{}.

\begin{quote}
Abends kommen die Mücken durch das offene Fenster.
\end{quote}

\begin{quote}
\emph{In the evening the mosquitoes come in through the open window.}
\end{quote}

Plural: \emph{die Mücken}. In Austria: \emph{die Gelse}.

\begin{cousinweb}
\emph{Mücke} is an inherited Germanic word; EtymDB lists English \emph{midge}, Dutch \emph{mug} and Swedish \emph{mygg} among its cousins.
\end{cousinweb}

\begin{grammarlens}[title={What you've built}]
One more word for telling, explaining and describing at B1.
\end{grammarlens}

\subsection*{Your turn}
\begin{itemize}
\raggedright
  \item \emph{Say it:} \emph{die Mücke}
  \item \emph{Say it:} \emph{die Mücke}, once more
  \item \emph{Say it:} a sentence of your own with \emph{die Mücke}
  \item \emph{Recall:} say the German for the bedbug, then the German for the cockroach, then say \emph{die Mücke} again
\end{itemize}

\begin{tcolorbox}[breakable,colback=teal!4,colframe=teal!35!black,title={Before you move on}]
What does die Mücke mean? (\textquotedblleft{}The mosquito, the gnat\textquotedblright{}.) Say it once more.
\end{tcolorbox}
\section[stechen]{stechen — to bite, to sting (of insects); to prick}

\label{lesson:GE-C430-stechen}



\begin{quote}
Before the new one: say the German for the cockroach, then the German for the mosquito.
\end{quote}

\subsection*{You'll want to know: stechen}
\textbf{stechen} — \textquotedblleft{}to bite, to sting (of insects); to prick\textquotedblright{}.

\begin{quote}
Eine Mücke hat mich am Arm gestochen.
\end{quote}

\begin{quote}
\emph{A mosquito bit me on the arm.}
\end{quote}

Strong verb: \emph{sticht}, \emph{stach}, \emph{hat gestochen}.

\begin{cousinweb}
\emph{stechen} is inherited through Old High German \emph{stehhan} from Proto-Germanic \emph{stikaną} 'to pierce' (EtymDB). English \emph{stick} is a cousin.
\end{cousinweb}

\begin{grammarlens}[title={What you've built}]
One more word for telling, explaining and describing at B1.
\end{grammarlens}

\subsection*{Your turn}
\begin{itemize}
\raggedright
  \item \emph{Say it:} \emph{stechen}
  \item \emph{Say it:} \emph{stechen}, once more
  \item \emph{Say it:} a sentence of your own with \emph{stechen}
  \item \emph{Recall:} say the German for the cockroach, then the German for the mosquito, then say \emph{stechen} again
\end{itemize}

\begin{tcolorbox}[breakable,colback=teal!4,colframe=teal!35!black,title={Before you move on}]
What does stechen mean? (\textquotedblleft{}To bite, to sting (of insects); to prick\textquotedblright{}.) Say it once more.
\end{tcolorbox}
\section[krabbeln]{krabbeln — to crawl (of insects, babies)}

\label{lesson:GE-C430-krabbeln}



\begin{quote}
Before the new one: say the German for the mosquito, then the German for to bite.
\end{quote}

\subsection*{You'll want to know: krabbeln}
\textbf{krabbeln} — \textquotedblleft{}to crawl (of insects, babies)\textquotedblright{}.

\begin{quote}
Eine Ameise krabbelt über das Kopfkissen.
\end{quote}

\begin{quote}
\emph{An ant crawls across the pillow.}
\end{quote}

Perfekt with \emph{sein} for moving: \emph{das Baby ist zur Tür gekrabbelt}.

\begin{cousinweb}
\emph{krabbeln} came into German from Middle Low German \emph{krabbelen} 'to crawl' (EtymDB), a cousin of Dutch \emph{krabben} 'to scratch'.
\end{cousinweb}

\begin{grammarlens}[title={What you've built}]
One more word for telling, explaining and describing at B1.
\end{grammarlens}

\subsection*{Your turn}
\begin{itemize}
\raggedright
  \item \emph{Say it:} \emph{krabbeln}
  \item \emph{Say it:} \emph{krabbeln}, once more
  \item \emph{Say it:} a sentence of your own with \emph{krabbeln}
  \item \emph{Recall:} say the German for the mosquito, then the German for to bite, then say \emph{krabbeln} again
\end{itemize}

\begin{tcolorbox}[breakable,colback=teal!4,colframe=teal!35!black,title={Before you move on}]
What does krabbeln mean? (\textquotedblleft{}To crawl (of insects, babies)\textquotedblright{}.) Say it once more.
\end{tcolorbox}
